% FILE: 02_lineage_and_context.tex
% PROJECT: ma — M&A Claim Taxonomy and Board Projection Surface

\section{Lineage and Context}
\label{sec:ma-lineage}

\subsection{2016 Language-Model Lesson}
\label{sec:ma-lineage-2016}

The thesis lineage opens in 2016 with a language-model experiment whose finding was not about
language models per se but about the fundamental inadequacy of local text imitation as a model
of consequential enterprise behavior. A system that predicts the next token in a sentence---or
the next step in a process description---does not conserve consequence. It can produce fluent
process narratives that describe processes which never ran, cannot be replayed, and contain
no evidence of actual organizational execution.

The M\&A domain makes this lesson acute. A PowerPoint deck presenting synergy projections is
precisely a language-model artifact: it is a fluent narrative about an organizational future
that may or may not be grounded in the operational reality of either the acquirer or the target.
The 2016 lesson---that imitation without process truth is dangerous---is, in the M\&A domain,
a capital-market failure mode. Buyers pay merger premiums for operational claims that are no
more grounded than next-token predictions.

The \texttt{ma/} module does not merely diagnose this failure. It defines the replacement
infrastructure: a claim taxonomy in which every assertion is grounded in an event log, every
event log is cryptographically chained, and every conformance result is independently replayable.

\subsection{Enterprise Process Gap}
\label{sec:ma-lineage-gap}

The 2016 lesson opens onto a structural enterprise process gap: the gap between what
organizations claim about their processes and what their event logs, if examined, would reveal.
This gap is not a minor measurement error. It is systematic. Management-prepared operational
presentations routinely omit rework intensity, compliance violations, SLA breach frequency,
and shadow-IT workarounds---not necessarily through deception but through the absence of any
methodology to measure and present these facts in a form an external party can verify.

The Operational Debt Taxonomy (source: \texttt{define\_operational\_debt\_taxonomy.md}) names
four classes of this gap:

\begin{itemize}
  \item \textbf{Process Spaghetti}: High trace entropy $H(L) > 3.0$ indicating variant
    explosion. Haircut impact: 10--15\% valuation discount.
  \item \textbf{Compliance Deficits}: Systematic rule violations exceeding five per 10,000
    cases. Haircut impact: 5--20\% valuation discount.
  \item \textbf{Legacy Lock-In}: Dependency on systems whose latency exceeds $3\times$ modern
    equivalents. Integration cost impact: 10--30\%.
  \item \textbf{Shadow IT Workarounds}: High ``manual update'' event rates indicating
    disconnected systems requiring undocumented human steps to function.
\end{itemize}

The Hostile Takeover Receipt taxonomy (source: \texttt{hostile-takeover-receipts.md}) observes
that incumbent management in a hostile scenario actively conceals these classes of operational
debt---and that HTR replay surfaces them precisely because replaying event streams without the
undocumented manual steps causes the replay to fail, making the ghost dependencies visible.

\subsection{Relationship to the Chatman Equation}
\label{sec:ma-lineage-chatman-equation}

The Chatman Equation $A = \mu(O^*)$ expresses the core thesis claim: an agent's authority $A$
over an outcome is a function $\mu$ of the operational star $O^*$---the complete, replayable,
cryptographically attested record of how outcomes were achieved. In the M\&A context, the
equation has a direct financial interpretation.

A seller who can produce $O^*$---a complete OCEL 2.0 event log with process model, conformance
certificate, and receipt chain---has demonstrable authority over the claimed operational outcome.
The merger premium is the market's valuation of $\mu(O^*)$: the premium for certainty over
estimation, for replayability over narration. A seller who cannot produce $O^*$ is selling
$\mu(\emptyset)$---a claim whose authority is zero because its operational basis is
unverifiable.

The Board Claim Taxonomy (source: \texttt{BOARD\_CLAIM\_TAXONOMY.md}) operationalizes this
through six tiers. T1 (Existence) and T2 (Coverage) claims correspond to $\mu(\emptyset)$:
they assert that a process exists or that metrics have been counted, but they are not buyer-
verifiable because they cannot be replayed. T3 (Conformance) through T6 (Recovery) claims
correspond to increasing levels of $O^*$ completeness: each requires an event log, a
process model, a conformance certificate, and---at T5 and T6---a live wasm4pm query or a
full variant-analysis and remediation receipt. The thesis argues that M\&A valuations should
price the tier of the seller's claims, not merely their numerical value.

\subsection{Relationship to Receipts and Replay}
\label{sec:ma-lineage-receipts}

The receipt infrastructure defined in the \texttt{ma/} module is not a new invention for M\&A.
It is the same receipt chain that \texttt{ggen} and \texttt{wasm4pm} use for operational
process evidence---extended to the capital-market surface. Each JSON receipt in
\texttt{receipts/} contains:

\begin{itemize}
  \item A \texttt{slide\_id} UUID linking the receipt to a specific board presentation slide.
  \item A \texttt{target\_log\_hash}: \texttt{BLAKE3} hash of the event log against which the
    claim was computed.
  \item A \texttt{process\_model\_hash}: \texttt{BLAKE3} hash of the Petri net or process tree
    used for conformance checking.
  \item A \texttt{query\_definition} referencing the \texttt{wasm4pm} engine and query URI.
  \item \texttt{verification\_results} with numeric fitness, precision, throughput, and
    financial impact values.
  \item A \texttt{validator\_signature}: Ed25519 signature over the RFC 8785 JCS-canonical
    form of the receipt JSON.
\end{itemize}

The receipt for EBITDA optimization (\texttt{rec\_ebitda\_rework\_001.json}) illustrates the
structure concretely: slide UUID \texttt{8a3e811c-...}, BLAKE3 log hash
\texttt{ee9ab723...}, fitness 0.982, precision 0.945, projected EBITDA impact \$1,250,000,
Ed25519 signature \texttt{ed25519\_sig\_8a3e811c\_02e0df2f...}. This is not narration. It is
a cryptographically signed commitment to a specific numerical claim derivable from a specific
event log by a specific query---independently executable by any party with the VDR contents.
